\section{Experimental design}
\label{sec:design}

\subsection{Data}
\label{sec:data}

Table~\ref{tab:datasets} summarises both corpora. The aviation corpus is
described in \citet{felendler2026trajrag}. The maritime corpus is new. We
downloaded 44 days of raw AIS position reports (19 April--1 June 2025, the
same calendar window as the aviation collection) from the NOAA Marine
Cadastre \citep{noaa2025marinecadastre}, a public United States government
archive with no access restrictions, which supports replication. Two
corridors with contrasting geometry were selected. The Houston Ship Channel,
from the Galveston Bay entrance to the Port of Houston, is a long, narrow
inland channel with dense tanker and cargo traffic. The Los Angeles/Long
Beach approach is a shorter, more open approach. From 26.2 million reports
we reconstructed 3,506 cargo and tanker voyages, segmented them into phases
and attached the marine environmental context. For testing we randomly
selected 250 vessels (125 per corridor) with at least four 5-minute points
in every phase, the same minimum segment length as in aviation; selection is
a uniform random sample over the full collection window, with no
chronological holdout by itself. AIS-ACNet's training and validation sets
exclude all test voyages and are themselves chronologically split
(Section~\ref{sec:baselines}).

An earlier implementation of TrajRAG's retrieval had two problems. First, it
excluded only the exact query voyage from a corridor's corpus-wide index,
not the other 124 test voyages, so a query could retrieve another test
voyage as a precedent -- retrieval memory contamination, not a disclosable
design choice. Second, the memory had no chronological cutoff, so a query
could retrieve a precedent recorded after it, which a real-time deployment
could never do. We rebuilt the memory to fix both: a corridor's retrieval
index is now built only from the chronologically earliest half of its
non-test voyages (1,141 of 2,283 non-test Houston voyages; 486 of 973
non-test Los Angeles/Long Beach voyages), and retrieval additionally
excludes every voyage in the test set, not only the query. Every TrajRAG
reconstruction reported in this paper queries only this pre-cutoff memory,
so no retrieved precedent can postdate the query, by construction rather
than by an exclusion list.

Reconstructing under the corrected memory, and under the corrected
preprocessing of Section~\ref{sec:preprocessing}, requires fresh LLM calls
for every query voyage, which this paper's budget could not extend to the
full 125-per-corridor test set together with the ablation and
matched-forecaster experiments of Sections~\ref{sec:maritime-ablations} and
\ref{sec:design} below. We therefore report TrajRAG's headline maritime
results on a smaller, still randomly-drawn common test set: 50 Houston and
10 Los Angeles/Long Beach voyages, sampled from the subset of the original
125-per-corridor test voyages that postdate each corridor's memory cutoff
(68 and 67 such voyages were available, respectively). Every other method
in Table~\ref{tab:extended} and Table~\ref{tab:corridor} -- Kalman, ARIMA,
LSTM, BiLSTM, Mean-Retrieval, Numeric-kNN and AIS-ACNet -- is rescored on
this identical 50/10-voyage set for a matched comparison; AIS-ACNet is
additionally reported on the full 125-voyage set as a robustness check that
the smaller matched sample is not anomalous (Section~\ref{sec:results}).
Section~\ref{sec:threats} discusses the resulting loss of statistical power,
particularly at Los Angeles/Long Beach.

\paragraph{Preprocessing correction}
\label{sec:preprocessing}
The 5-minute resampling grid originally interpolated latitude and longitude
with a time-weighted method that could use an AIS report recorded
\emph{after} a given grid timestamp, including the single most load-bearing
point in the whole pipeline: the last point of the verified prefix that
every method treats as ground truth to reconstruct from. Every other
resampled field (SOG, COG, heading, \ldots) already used a causal
last-observation-carried-forward rule; we moved latitude and longitude onto
the same rule, removing the leak. Rescoring the deterministic classical and
numeric baselines (Kalman, ARIMA, Numeric-kNN) on the corrected data versus
the original shows this is not a negligible effect: every one of 12
method/phase/corridor cells we checked became worse under the correction,
by 2--40\%, confirming the leak had been making reconstruction look easier
than it is. All results in this paper use the corrected data.

\begin{table}[t]
\centering
\caption{Datasets. The aviation corpus is the one used in the conference
version \citep{felendler2026trajrag}; the maritime corpus is new.}
\label{tab:datasets}
\footnotesize
\begin{tabular}{p{2.5cm}p{4.3cm}p{4.3cm}}
\toprule
 & Aviation & Maritime \\
\midrule
Source & FlightRadar24 ADS-B & NOAA Marine Cadastre AIS (public) \citep{noaa2025marinecadastre} \\
Corridors & Paris$\to$Rome; Barcelona$\to$Munich & Houston Ship Channel; Los Angeles/Long Beach approach \\
Collection window & 19 Apr--1 Jun 2025; 17 Dec 2025--18 Jan 2026 & 19 Apr--1 Jun 2025 (44 days) \\
Raw volume & 46,550 messages (Paris$\to$Rome) & 26.2 million position reports \\
Trajectories & 1,544 flights (931 + 613) & 3,506 voyages (2,408 Houston; 1,098 Los Angeles/Long Beach) \\
Sampling & $\approx$120\,s (native) & 300\,s (uniform resampling) \\
Test set & 100 flights, chronologically held out (Paris$\to$Rome) & 50 Houston + 10 Los Angeles/Long Beach vessels, each with $\geq 4$ points per phase, drawn from a 250-vessel (125 per corridor) pool and postdating the retrieval memory's chronological cutoff \\
Environmental context & Meteomatics & Open-Meteo Marine \citep{openmeteo2025} \\
\bottomrule
\end{tabular}
\end{table}

\subsection{Blackout protocol}

We simulate post-detection blackouts on real, uncorrupted trajectories, so
that the true positions are available as ground truth. Following a
safety-oriented, attack-agnostic stance, we model every attack (gradual
drift, velocity manipulation or capture-and-pull-off) the same way: all
satellite-derived fields after onset are withheld. In aviation, onset times
are randomised and blackouts last 5--15 steps (10--30\,min)
\citep{felendler2026trajrag}. In the maritime benchmark, each phase segment
of $N$ points is split at $\lfloor N/2 \rfloor$. The first half is the
verified prefix, and the second half is withheld and must be reconstructed.
This gives mean blackout durations of 37.2, 35.2 and 66.1\,min for
departure, transit and arrival on the Houston test set. The maritime
protocol is applied identically to every maritime method, including every
baseline: no baseline receives any field reported after the onset index.
Because the two domains use different protocols and sampling intervals,
absolute errors are compared within a domain, not across domains.

\subsection{Baselines}
\label{sec:baselines}

\textbf{Classical, recurrent and retrieval-only.} These are a Kalman filter
(a constant-velocity model that accounts for currents in the maritime case,
using only prefix data and environmental context), ARIMA fitted to each
verified prefix, LSTM and BiLSTM networks trained on the non-test voyages of
each phase (about 1,600--3,300 per phase), and Mean-Retrieval, which
averages the continuations of the top-$K$ retrieved segments without an
LLM.

\textbf{Numeric-kNN.} This baseline is a non-parametric retrieval model
operating on numbers rather than text. It finds the five training prefixes
of the same phase closest to the test prefix by mean point-wise Haversine
distance and predicts each withheld step as the coordinate-wise mean of
their continuations. It isolates the value of textual states combined with
LLM blending over plain numeric nearest neighbours.

\textbf{AIS-ACNet.} We retrain AIS-ACNet \citep{shin2024aisacnet} from its
public implementation on our voyages (on the corrected data of
Section~\ref{sec:preprocessing}), using a chronological 85/15
train/validation split of the non-test Houston voyages, and evaluate it on
the common test set with the same blackout boundary. We report it both on
the full 125-voyage test set and on the 50/10-voyage set TrajRAG is
evaluated on, to check that the smaller matched sample is not anomalous
(Section~\ref{sec:results}). AIS-ACNet's
convolutional receptive field is fixed, so we train one model per phase,
with the window set near the 95th percentile of the true half-phase length
(Table~\ref{tab:acnet-config}). Voyages whose true withheld half exceeds the
window are truncated to it for both the model's output and the scoring mask
-- 4, 6 and 5 of the 125 test voyages in departure, transit and arrival
respectively (3--5\% per phase); shorter ones are padded with edge values
and masked out of the loss and the metric. For the remaining 95--97\% of
voyages per phase, AIS-ACNet reconstructs exactly the same horizon as
TrajRAG; for the truncated tail, AIS-ACNet is scored only on the first
\emph{window} steps of the blackout, not the full one TrajRAG reconstructs.

\begin{table}[t]
\centering
\caption{Phase-specific AIS-ACNet configuration. The window is set near the
95th percentile of the true half-phase length so that most voyages are
predicted over their full withheld half.}
\label{tab:acnet-config}
\footnotesize
\begin{tabular}{lccc}
\toprule
Phase & Window (steps) & Residual blocks & Test voyages scored \\
\midrule
Departure & 21 & 7 & 125 \\
Transit & 15 & 5 & 125 \\
Arrival & 32 & 11 & 125 \\
\bottomrule
\end{tabular}
\end{table}

\textbf{Modern forecasters.} DLinear, PatchTST, iTransformer, Informer,
Autoformer, FEDformer and TimesNet are trained with the Time-Series-Library
implementations \citep{wu2023tslib} on eight features per step (latitude,
longitude, SOG, COG, heading, wave height, 10\,m wind speed and current
speed). The features are z-scored with training-split statistics.
Encoder--decoder models receive the tail of the \emph{observed} window as
their decoder warm-up, so no target value leaks into the input. These
architectures need fixed input and output lengths. Following the aviation
harness of the conference paper, we use a small exact-fit window per phase
(4, 3 and 4 steps for departure, transit and arrival) and include a voyage
only if its half-phase length matches the window. This retains about 80\%,
71\% and 77\% of eligible voyages. The forecasters therefore solve an
easier, shorter-horizon problem on a subset of voyages, and
Table~\ref{tab:extended} reports them in a separate block.

\subsection{Metrics and statistical testing}
\label{sec:metrics}

The primary metric is the mean Haversine error (MHE) \citep{zhao2024space}
over the withheld steps $T$,
\begin{equation}
\mathrm{MHE} = \frac{1}{|T|} \sum_{i \in T} d_{\mathrm{hav}}(\hat{y}_i, y_i),
\label{eq:mhe}
\end{equation}
reported per phase. A given absolute error means more for a slow ship than
for a fast one, so we also normalise MHE by two quantities that are
independent of the method and computed from ground truth: the mean distance
travelled during the blackout and the mean blackout duration. Differences
between TrajRAG and each full-horizon baseline are tested with paired
two-sided Wilcoxon signed-rank tests \citep{wilcoxon1945individual} on
per-voyage errors of the same voyages (50 Houston, 10 Los Angeles/Long
Beach). We report raw $p$-values and note
which conclusions survive a Bonferroni correction. We additionally report,
for each comparison, the median paired difference (baseline minus TrajRAG)
with a 95\% bootstrap confidence interval over 10,000 resamples of voyages,
and the median, 90th percentile and maximum of the per-voyage error
distribution for every full-horizon method.

\subsection{Implementation}
\label{sec:implementation}

Both instantiations use \texttt{text-embedding-3-large} as $g$ and FAISS
inner-product indices \citep{johnson2021billion}. The aviation adapter uses
\texttt{gpt-5.1} with $K=5$, the configuration selected in the conference
study. The maritime adapter uses \texttt{gpt-4.1} with $K=6$, fixed before
any maritime test-set result was observed; neither backbone was fine-tuned,
and no maritime hyperparameter was tuned on the test set. Per their
respective model documentation, \texttt{gpt-4.1}'s training data cutoff is
June 1, 2024 and \texttt{gpt-5.1}'s is September 30, 2024 -- both postdate
the Houston and LA/Long Beach voyage windows used in this study (19
April--1 June 2025 collection, but reflecting vessel traffic and port
operations that are not novel information the models could have memorised
verbatim). AIS-ACNet and the
seven forecasters are trained with Adam and early stopping (AIS-ACNet:
masked mean-squared error on position and auxiliary SOG/COG, patience 20, at
most 150 epochs; forecasters: mean-squared error, learning rate $10^{-3}$,
batch size 32, patience 15, at most 60 epochs). LLM decoding uses
temperature 0 for position generation and temperature 0.3 for the one-off
textual summary at blackout onset, in both domains.
